\clearpage
\subsubsection{Section Reproduction Steps}\label{subsubsec:reproduce-contacts}
From the repository root, after \texttt{make setup}, run:
\begin{lstlisting}[basicstyle=\ttfamily\small,breaklines=true]
make reproduce-contact-section
make report
\end{lstlisting}
The target reuses fingerprint-checked May--July 2026 inputs and the October 28,
2026 observation cutoff; no warehouse refresh occurs. It runs:
\begin{itemize}[leftmargin=*]
\item \path{opportunity-customer-contacts-by-cancel-reason} for
Table~\ref{tab:customer-contacts-by-cancellation-reason}; matching inbound
contacts have no creation-relative day bounds and end at the observed order endpoint.
\item \path{opportunity-aged-customer-contacts} for
Tables~\ref{tab:aged-customer-contacts-by-cancellation-day} and
\ref{tab:aged-customer-contact-day-frequency}, using elapsed cancellation days
$-13$ through $+3$ and all handled states.
\item \path{opportunity-contact-frequency-reorders} for
Table~\ref{tab:contact-frequency-same-pet-part-reorders} and
\path{opportunity-contact-window-reorders} for
Table~\ref{tab:contact-window-same-pet-part-reorders}. Both use successful
same-customer, same-pet and same-part-number reorders within 14 days.
\end{itemize}
The last target also regenerates the frequency study from its verified cache
before calculating the contact-window comparison. Retained inputs under
\path{data/opportunity/} are required; they are not in Git. For new Snowflake
snapshots, run the underlying extraction targets without \texttt{CACHED=1}
using the same dates, then regenerate the comparisons. See HOWTO.md.
